% Source: https://artofproblemsolving.com/community/c3h1662953_mathcounts_warmup_help
% Stable identity: aops-topic-1662953-q01
Since Du passes Priya, he is faster. Two runners starting together are side by side again for the first time when the faster one has run exactly one lap more than the slower one. When Du finishes his $4$th lap he is passing Priya for the first time, so at that moment Priya has run exactly $4-1=3$ laps. Their speeds are therefore in the ratio $4:3$; say Du runs $4$ laps and Priya $3$ laps per unit of time. Amanda then runs $\dfrac{4+3}{2}=\dfrac{7}{2}$ laps per unit of time.

Amanda gains $\dfrac{7}{2}-3=\dfrac{1}{2}$ lap per unit of time on Priya, so she first passes her when she is one lap ahead, after $2$ units of time. By then Amanda has completed
\[
2\cdot\dfrac{7}{2}=\boxed{7}\text{ laps}
\]
(and Priya $6$).
